% docs/style/dfd-macros.tex -- local copy of tex-skripten/formatierungen/dfd-macros.tex
% Copied 2026-10-12 so that the repository compiles stand-alone (GitHub).
% Local version: 1.0.0 -- unchanged copy

% dfd-macros.tex
% ==========================
% DFD Diagram Makros
% ==========================
\usepackage{ifthen}


\usetikzlibrary{shapes.geometric, arrows.meta, positioning}

% 1. Komplette Stil-Definitionen
\tikzset{
    prozess/.style = {
        circle, draw=black, thick, fill=blue!20,
        minimum size=1.8cm, align=center
    },
    io/.style={
        rectangle, draw=black, thick, fill=green!20,
        minimum height=1.2cm, minimum width=2cm, align=center
    },
    storage/.style={
        inner ysep=3pt,
        append after command={
            \pgfextra{\draw (\tikzlastnode.north west) -- (\tikzlastnode.north east);
                \draw (\tikzlastnode.south west) -- (\tikzlastnode.south east);}
        }
    },
    dataserver/.style = {
        cylinder, shape border rotate=90, draw, thick, fill=gray!20,
        minimum height=1.5cm, aspect=0.4, text width=2.5cm, align=center
    },
    guinode/.style = {
        rectangle, rounded corners, draw, thick, fill=orange!20,
        minimum height=1.2cm, minimum width=2.5cm, align=center
    },
    external/.style = {
        rectangle, draw, thick, fill=gray!30,
        minimum height=1.2cm, text width=2.8cm, align=center
    },
    % KORREKTUR: Stil umbenannt
    ausgabe/.style = {
        trapezium, trapezium left angle=70, trapezium right angle=110,
        draw, thick, fill=red!10, text width=2.5cm, align=center
    }
}

% 2. Makros fuer jeden Stil
\newcommand{\prozess}[3][]{\node[prozess, #1] (#2) {#3};}
\newcommand{\io}[3][]{\node[io, #1] (#2) {#3};}
\newcommand{\storage}[3][]{\node[storage, #1] (#2) {#3};}
\newcommand{\dataserver}[3][]{\node[dataserver, #1] (#2) {#3};}
\newcommand{\guinode}[3][]{\node[guinode, #1] (#2) {#3};}
\newcommand{\external}[3][]{\node[external, #1] (#2) {#3};}
\newcommand{\ausgabe}[3][]{\node[ausgabe, #1] (#2) {#3};}

% ===== DOKUMENTEN-EINSTELLUNGEN =====
%\geometry{a4paper, top=2.5cm, bottom=2.5cm, left=2.5cm, right=2.5cm}
